\section{毕业设计（论文）内容概述}

\subsection{项目来源及开发目的和意义}

本课题来源于面向蛋白质设计任务的多 Agent 系统工程实践。现有系统已经具备显式有限状态机（FSM）、四类 Agent 职责边界、人工介入（HITL）、分层恢复链和工具适配等基础能力，能够围绕蛋白质设计任务形成较完整的工程闭环。但若仅将其理解为“生命科学助手系统”，论文的重点就容易停留在任务分解、工具接入和结果展示层面，难以形成边界清晰、能够独立成立的核心算法问题。

造成这一问题的根本原因，并不在于系统功能不足，而在于研究对象此前定义得过于贴近应用场景。若始终将课题表述为“面向蛋白质设计的多 Agent 辅助系统”，那么论文的重点就容易落在任务分解、工具接入、结果展示和系统工程组织上；即使加入门禁、恢复和人工决策机制，也仍然容易被视为生命科学助手系统的增强版本，而难以形成一个边界清晰、可被论证的计算机专业研究主线。

基于这一判断，本文将课题界定为面向高代价蛋白质设计工作流的自适应工具链规划与运行时控制问题。在这一问题定义下，蛋白质设计是验证算法有效性的应用场景；系统现有的 FSM、HITL、恢复链、日志与快照等能力，则构成核心算法可依附的运行环境。论文要回答的核心问题因此转变为：在一个高代价、长链路、可失败、可恢复、允许人工介入的科学工作流中，系统应依据什么原则决定下一步继续执行、局部修补、后缀重规划或终止止损。

围绕这一问题，本文在现有多 Agent 系统基础上，进一步设计并逐步实现了面向高代价蛋白质设计工作流的自适应工具链规划算法。依据当前设计文档中的核心算法规范，该算法并不把问题简化为“从若干工具中选一个工具”，而是把研究重点放在工作流级候选链生成、静态多目标评分、运行时状态更新、运行时重排序和动作选择上：一方面，系统在规划、修补和重规划阶段生成 Top-K 结构化候选，并综合可执行性、目标贴合度、风险、成本和恢复复杂度进行排序；另一方面，系统在执行过程中维护轻量运行时状态，对当前链路的成功概率、结构性失败倾向、恢复余量、剩余成本和证据充分性进行持续估计，并据此在 \texttt{continue}、\texttt{patch\_local}、\texttt{suffix\_replan} 和 \texttt{stop} 等治理动作之间做出选择。

上述转变使本课题的意义发生了实质性提升。首先，从研究层面看，本文不再只是构建一个多 Agent 蛋白质设计平台，而是尝试将高代价科学工作流中的执行治理、恢复控制与人机协同抽象为一个部分可观测的序贯决策问题。其次，从工程层面看，本文所提出的方法并非脱离现有系统另起炉灶，而是在现有 FSM、Agent 边界和审计链上增量融入，使算法设计与系统落地保持一致。最后，从论文表达层面看，围绕“自适应工具链规划与运行时控制”这一主线，原本分散在架构、恢复、门禁和等待态中的多个概念得以被统一解释为同一核心算法问题的不同组成部分，这也为后续的实验设计和结果分析提供了更清晰的结构基础。

\subsection{主要开发任务}

围绕新的研究主线，中期阶段的主要开发任务可概括为四个方面。

第一，完成研究问题的重构与核心算法定义。具体包括将原先偏应用系统的表述，收束为高代价科学工作流中的控制与规划问题；明确算法的输入、输出、优化目标和边界约束；并将其落实为“静态候选规划、运行时状态更新与重排序、运行时动作选择”三层结构。

第二，完成算法与现有系统运行骨架的对齐。该部分工作要求算法不能破坏已有 FSM、HITL 和四类 Agent 职责边界，而是应以增量方式嵌入当前工作流：Planner 继续负责候选生成与静态评分，Workflow Runtime 负责状态推进与动作生效，Executor 负责执行与失败检测，Safety 和 Summarizer 则分别提供风险观测和结果汇总，从而保证“算法增强”不会退化为“系统重写”。

第三，完成运行时状态、动作选择和审计证据的落地。除初始阶段已经形成的候选门禁、人工确认和恢复链外，后续迭代重点补充了运行时状态契约、状态快照恢复、等待态运行时摘要、动作选择器、影子重排序接口以及面向对照实验的观测基线等内容，使“控制依据”能够从原先偏静态的启发式逻辑，逐步转化为结构化、可追踪、可比较的运行时决策过程。

第四，完成论文材料、实验材料与展示材料的同步整理。该部分工作要求将设计文档、代码实现、issue 迭代、验证报告和图示材料统一转写为围绕核心算法展开的论证链条，并为横向对比实验和工作量统计补充保留清晰接口。

\subsection{本人所承担任务（模块）说明}

在本课题实施过程中，本人承担的工作主要集中在研究问题收束、核心算法设计、系统实现推进和论文材料整理四个方面。

首先，在研究定义层面，本人持续梳理系统现状、已有不足及选题风险，推动课题从“生命科学多 Agent 助手系统”转向“面向高代价科学工作流的自适应工具链规划与运行时控制问题”，并据此补充算法定义、创新边界和论文主线。

其次，在算法与系统实现层面，本人重点参与了候选工具链生成、静态多目标评分、运行时状态契约、动作选择控制流、恢复链审计以及等待态快照恢复等关键内容的设计与落地。这部分工作并非单纯增加工具或补充界面，而是围绕核心方法逐步建立起从任务输入、候选生成、运行时观测、状态更新到动作生效的统一控制链。

再次，在系统工程推进层面，本人参与完成了 FSM/HITL 主链、六阶段蛋白质设计工作流、工具接入、远端服务联动、验证报告沉淀和实验运行器扩展等工作，使新算法具备能够在真实系统中运行和被验证的基础环境。随着项目迭代推进，相关工作进一步聚焦到运行时状态、动作选择和横向实验运行基础的补齐。

最后，在论文与答辩材料方面，本人负责将设计文档、实现进展、issue 历史、验证材料和图示资源整理为可直接服务于中期论文和答辩展示的书面内容。关于项目工作量的定量统计，如代码规模、阶段性交付数量、issue/PR 数量及测试覆盖数据等，本文将在统一统计后于相应位置补充具体数值。
